Row scaling preserves a mixed-integer programming model in exact arithmetic but
changes the application-unit meaning of absolute feasibility tolerances. We
develop a model-generation workflow using original-unit residual budgets. It
selects row factors under coefficient limits, reports incompatible requirements,
accounts for integer rounding, checks binary64 export errors during selection,
and independently checks returned points. Evaluation combines 306 hydrothermal
dispatch models, kernel-matched and Ruiz controls, enumerated stress tests on
three solvers, and independent batch scheduling. In prospective validation on
24 perturbed scenarios from six known inputs, the budget-aware Euclidean policy
has paired geometric-mean work and complete-process time ratios of 0.422 and
0.797 against no external scaling. All policies meet the physical budgets;
one budget-aware bound is independently contradicted, and candidate cost
deterioration reaches 1.02\%. A larger fixed-policy study comprises 3,840 runs
on 160 models from 40 historically generated parents and two solvers. Against
conventional Euclidean scaling, paired ratios of penalized verified-availability
time are 0.995 for Gurobi and 1.000 for CPLEX, with whole-parent uncertainty
intervals containing one. Gurobi complete-process time increases by 7.9\%.
Every eligible final point meets the mapped physical budgets under every policy;
strict written-bound checks limit CPLEX's combined acceptance. These within-family
results support explicit precision specification, admissible selection and
reproducible verification. They establish neither general incremental acceleration,
additional physical-budget recovery in the prospective samples, nor certified
mixed-integer optimality.
